\begin{theorem}[Proved: cyclic four-object support through seventeen]
\label{CP2:A-strong-four-seventeen}
A four-object category with cyclic support and at most
seventeen arrows cannot have type $\CP^2$.
\end{theorem}
\begin{proof}
Suppose $BC\simeq\CP^2$. In particular the category is
connected and $\dim H^2(BC;\mathbb Q)=1$, as required
by the binary-stalk obstruction below. If an object in a proper strong
component is a retract of another, its deletion removes
at least four arrows: its identity, arrows to and from
another object of the component, and an external incident
arrow. The last exists because every object of a strong
component inherits any external support adjacency. If the
whole category is strong, deletion removes at least seven
arrows. Thus an object retract always reduces to a
three-object category with at most thirteen arrows,
excluded by Theorem~\ref{CP2:A-high-all-eight}. Assume
there is no object retract.

A strong component on $r$ objects has its $r$ identities
outside its categorical minimal ideal, so costs at least
$r^2+r$ arrows. A wholly strong four-object category
therefore costs at least twenty. A proper three-object
strong component costs at least twelve and has at least
three interface arrows to the remaining corner. That
corner has order at most two. The strong component has
finite-group type: otherwise its total row and column
counts would both be at least four and its internal cost
would already be nineteen. Both blocks have zero second
rational cohomology, so the two-block lemma applies.
If there are two strong pairs, each costs at least six
and their interface at least four. Each internal pair
therefore has at most seven arrows and finite-group type.
The same two-block argument excludes this case.

There remains one strong pair $S$ and two singleton
components. Put $q=|S|$. The outside components have
corner monoids of orders $b,c$.

\emph{Extreme pair.} By opposition suppose $S$ is a
source. Let $D$ be its two-object complement, and let
$h$ count the interface arrows. Connectedness gives
$q\le12$. If $q=11$ or $12$ but $H^2(BS;\mathbb Q)=0$,
then $|D|\le4$ since $h\ge2$; the second-homology cost
lemma gives $H^2(BD;\mathbb Q)=0$, and the two-block
lemma applies. For $q\le10$, the row-column bound makes
$BS$ of finite-group type, so $H^2(BS;\mathbb Q)=0$.
If also $H^2(BD;\mathbb Q)=0$, the two-block lemma
finishes. If the latter group is nonzero, then $|D|\ge6$:
for connected $D$ this is the second-homology cost lemma,
and for disconnected $D$ one corner has order at least
five and the other at least one. If $q=6$ or $7$, the
minimal group is trivial and each object has one row
and one column. The actual source-component orbit
collapse saves at least $q$ arrows, as proved in
Theorem~\ref{CP2:A-high-all-eight}, reducing to the
three-object bound. The case $q=10$ is impossible here,
since $q+|D|+h\ge10+6+2>17$.

It remains in this branch that $q=8$ or $9$. The complement
must be connected: otherwise it costs at least six and
has at least four interface arrows, exceeding seventeen.
Moreover $S$ can meet only one of its objects, say $u$,
since meeting both would again require four interface
arrows. The outside edge is directed $v\to u$; an edge
in the other direction would force another interface.
The inequality is now $q+|D|+h\le17$, with $|D|\ge6$
and $h\ge2$. At $q=9$ this forces $|D|=6,h=2$; at
$q=8$ either $h=2$, or $|D|=6,h=3$.

The minimal group is trivial and one total row or column
count is two. For $h=2$, all external stalks are singleton,
so the one-row source collapse or the fixed-exit column
dual collapses the pair, saving at least seven arrows.
For $q=8,h=3$, a row total of two gives the source
collapse immediately. Otherwise the totals are three
rows and two columns, and $S$ consists exactly of its
six-element minimal band and two identities. At its
one-row object all incoming internal arrows are fixed
by the local minimal idempotent, and its only unit is
its identity. There are no external entrances. The
one-sided object-deletion lemma applies, removing at
least its two corner arrows, three internal cross arrows,
and one external arrow. This again leaves at most
thirteen arrows on three objects.

Now suppose $H^2(BS;\mathbb Q)\ne0$. Then $q=11$ or
$12$, with trivial minimal group and total rows and
columns both three. If $D$ is disconnected, the only
budget is $q=11$, two trivial outside corners, and four
singleton interface stalks. This is the full join of
$S$ with two discrete objects, hence has zero products.
If $D$ is connected, $S$ meets only its final object $u$,
as before. We have $|D|+h\le6$. If $h=2$, then
$|D|\le4$ and the corner of $u$ has order at most two.
The full category on $S\cup\{u\}$ is the join of $S$
and that corner and is rationally acyclic. Its nerve
and $ND$ cover the whole nerve, intersecting at the
rationally acyclic corner $u$; $H^2(BD;\mathbb Q)=0$
by its size. Hence the whole second cohomology is zero.
The only remaining case is $q=11$, $|D|=3$, $h=3$.
Here $D$ is a single interval between trivial corners,
so gluing it at $u$ does not change the homotopy type
of $S\cup\{u\}$. Each local row count of $S$ is at most
two and $u$ is trivial. The binary projective-cover
lemma kills every second product on this act cone.

\emph{Nonextreme pair.} The condensation is the full
chain $u<S<v$. Write $\alpha_i=|C(u,i)|$ and
$\beta_i=|C(i,v)|$ for its two entrance and exit
stalks, and $z=|C(u,v)|$. They are all positive. Put
\[
 P=\sum_i(\alpha_i-1)+\sum_i(\beta_i-1),\qquad
 \delta=(b-1)+(c-1)+(z-1).
\]
Then
\[
 q+7+P+\delta\le17,
 \qquad q\le10,\qquad P+\delta\le10-q.
\]
A nontrivial minimal group forces $q=10$, all five
external hom-sets singleton, and both outside corners
trivial. The initial object $u$ is then initial in
the whole category. Assume henceforth the minimal
group is trivial.

For $q\ge7$ and $z\le2$, both endpoint corners have
order at most four and $P\le3$. If both entrance
and exit sides had a stalk of size at least three,
they would cost at least four units of $P$. Thus at
least one endpoint has its two stalks at $S$ of size
at most two; its direct stalk $z$ is also binary.
Lemma~\ref{CP2:A-high-small-binary}, whose opposite
block may be any finite category, kills its square.

If $q\ge7$ and $z\ge3$, only $q=7$ or $8$ is possible.
For $q=8$ one has $P=0$; one row or one column at every
object and the fixed singleton interfaces permit the
block collapse. For $q=7$, $P\le1$. The pair consists
of its four-arrow minimal band, two identities, and
one extra arrow. Its local minimal idempotents are
zeros. Of the four conditions that one local idempotent
fix all incoming or all outgoing cross arrows, the
extra internal arrow can violate at most two, and
the external nonsingleton stalk at most one. At least
one condition holds. The zero guarantees the unit
condition in the one-sided deletion lemma. Deletion
removes at least six arrows, including one entrance
and one exit, so the three-object bound applies.

Finally let $q=6$. The pair is exactly its four-arrow
minimal band plus its two identities. All internal
cross arrows are fixed by the local zeros, so a failed
one-sided condition requires a nonsingleton external
stalk at that specific object and side. If an endpoint
has order five, the budget forces $P=0,z=1$ and the
pair collapses. Otherwise both endpoint corners have
order at most four. For $z\ge3$ one has $P\le2$,
so fewer than four fixed conditions can fail and
one-sided deletion applies. For $z\le2$, failure
of both small-endpoint binary tests requires a stalk
of size at least three on each side, using $P\ge4$.
Equality is forced, with $b=c=z=1$; each side has
one three-element stalk and one singleton stalk.
Thus two of the four fixed conditions still hold,
and one-sided deletion again applies. Every alternative
is excluded, completing the cyclic four-object scope.
\end{proof}
